% CLINT chapter. Configuration-driven: hart counts, vector numbers and
% addresses come from include/defines.tex. Section order is the per-peripheral
% template enforced by python/check_intro_names.py.
\section{Overview}

The Core-Local Interruptor (\peripheral{CLINT}) provides the two interrupt sources on which multi-core software is built: a per-hart software interrupt, used as the inter-processor interrupt (IPI), and a per-hart timer interrupt driven by one shared free-running 64-bit counter. The block sits in the shared window behind the multi-core arbiter, so any hart can raise, clear or program any hart's interrupts.

\begin{itemize}
	\item One software-interrupt register per hart (\register{MSIP0} to \register{MSIP\MaxHartIndex}), writable by any hart
	\item One shared free-running 64-bit counter (\register{MTIMEL}/\register{MTIMEH}), incrementing once per MCLK cycle
	\item One 64-bit compare register per hart (\register{MTIMECMP0L}/\register{MTIMECMP0H} to \register{MTIMECMP\MaxHartIndex L}/\register{MTIMECMP\MaxHartIndex H})
	\item Level-sensitive delivery on two hardwired vectors per hart: \ClintMsipVector{} (software) and \ClintMtipVector{} (timer)
	\item Wakes a hart parked by the \asminline{sleep} instruction, so parked harts consume no dynamic power
	\item Compare registers reset to all ones, so no timer interrupt fires before one is programmed
\end{itemize}

\section{Functional description}

\subsection{Software interrupts}

Writing 1 to \register{MSIPh} asserts the software-interrupt level into hart $h$; writing 0 deasserts it. The interrupt is a level, so the handler on hart $h$ must write 0 to its own \register{MSIPh} before returning, because the level would otherwise re-enter the handler immediately.

The software interrupt is the wake mechanism for a parked hart: a hart that has enabled vector \ClintMsipVector{} and executed \asminline{sleep} resumes when any hart writes 1 to its \register{MSIPh}. The woken hart's handler clears the level and executes \asminline{wake} before returning, because a handler that returns without \asminline{wake} puts the hart back to sleep (Section \ref{ss:sleep}). The boot ROM uses this same path to launch a freshly loaded tile (Section \ref{ss:multicore-boot}).

\subsection{Timer interrupts}

The counter is one 64-bit value shared by all \NumHartsWord{} harts and clocked by the free-running MCLK. Each hart owns one 64-bit compare value; the timer level into hart $h$ is asserted while the counter is greater than or equal to hart $h$'s compare value, and it is deasserted by moving the compare value into the future or back to all ones (the disarmed reset value).

The registers are 32 bits wide, so every 64-bit value is a low/high pair. The counter cannot be read atomically; Section \ref{ss:clint-programming} gives the read sequence that tolerates a carry between the two halves.

\subsection{Clocking}

The \peripheral{CLINT} runs on MCLK and is never clock-gated, because a hart whose own clock is stopped by \asminline{sleep} can only be woken by logic that keeps running. The counter therefore advances through every sleep state, and its rate follows the MCLK divider in \register{CLKDIVCR}.

\subsection{Address decoding}

The block occupies shared-window page 1 from \texttt{\ClintBaseAddress}. \register{MSIPh} is at \texttt{\ClintBaseAddress}${}+4h$, the counter pair starts at \texttt{\ClintMtimeAddress}, and hart $h$'s compare pair is at \texttt{\ClintMtimecmpBaseAddress}${}+8h$. Only the low address bits are decoded, so the registers alias every \ClintAliasBytes{} bytes up to \texttt{0x5FFF}; software uses the canonical addresses because the alias pattern is not part of the register interface and may change with the hart count.

\section{Interrupts and events}

Both vectors are hardwired into every hart's interrupt controller and do not pass through the \peripheral{IRQROUTER}: each hart receives only its own \register{MSIPh} and its own compare match. Table \ref{t:clint-irq} lists them.

\begin{table}[htb]
	\centering
	\caption{\peripheral{CLINT} interrupt vectors.}
	\label{t:clint-irq}
	\begin{tabular}{ l p{4.2cm} p{4.0cm} p{4.6cm} }
	\textbf{Vector} & \textbf{Condition} & \textbf{Set by} & \textbf{Cleared by} \\ \hline
	\ClintMsipVector{} & Software interrupt to hart $h$ & Any hart writing 1 to \register{MSIPh} & Writing 0 to \register{MSIPh} \\
	\rowcolor{tablehighlightcolor}
	\ClintMtipVector{} & Timer match for hart $h$ & Counter reaching hart $h$'s compare value & Writing a compare value above the counter \\
	\hline
	\end{tabular}
\end{table}

\section{Programming} \label{ss:clint-programming}

\subsection{Reading the counter}

\begin{enumerate}
	\item Read \register{MTIMEH}.
	\item Read \register{MTIMEL}.
	\item Read \register{MTIMEH} again; if it differs from the first read, repeat from step 1, because the low half wrapped between the two reads.
\end{enumerate}

\subsection{Arming a timer interrupt for hart $h$}

\begin{enumerate}
	\item Write all ones to \register{MTIMECMPhH}, so that no intermediate compare value can be below the counter.
	\item Write the low half of the target time to \register{MTIMECMPhL}.
	\item Write the high half of the target time to \register{MTIMECMPhH}.
	\item Enable vector \ClintMtipVector{} in hart $h$'s interrupt controller.
\end{enumerate}

\subsection{Sending a software interrupt to hart $h$}

\begin{enumerate}
	\item Write 1 to \register{MSIPh}.
\end{enumerate}

\section{Usage notes}

\paragraph{Clear the software-interrupt level before returning.} The level stays asserted until \register{MSIPh} is written 0; a handler that returns first is re-entered at once.

\paragraph{Execute \asminline{wake} in the handler of a sleeping hart.} A hart woken by an interrupt returns to sleep after the handler unless the handler executed \asminline{wake} (Section \ref{ss:sleep}).

\paragraph{Program the compare pair high half first.} Writing the low half first can momentarily produce a compare value below the counter and raise a spurious timer interrupt.

\paragraph{Do not send a software interrupt to a hart whose loader mailbox row is incomplete.} The boot ROM consumes the row when the parked hart wakes, so the row must be written before \register{MSIPh} (Section \ref{ss:multicore-boot}).
